\documentclass[11pt]{article}

\usepackage[margin=1in]{geometry}
\usepackage{amsmath,amssymb}

\newcommand{\R}{\mathbb{R}}

\begin{document}

\section*{Deliverable 1: Single-Segment Trajectory Optimization}

Consider the minimum-velocity problem
\begin{equation}
    \min_{P(t)} J(P)
    = \int_0^1 \left(P^{(1)}(t)\right)^2\,dt
    \quad \text{subject to} \quad
    P(0)=0, \qquad P(1)=1.
\end{equation}
Because $P(0)=0$, the constant coefficient of $P(t)$ is zero.

\subsection*{1. Polynomial of degree 1}

Let
\begin{equation}
    P(t)=p_1t.
\end{equation}
The terminal constraint gives
\begin{equation}
    P(1)=p_1=1.
\end{equation}
Therefore, $P'(t)=1$, and
\begin{equation}
    J(P)=\int_0^1 1^2\,dt=1.
\end{equation}
Hence,
\begin{equation}
    \boxed{P^*(t)=t, \qquad J^*=1.}
\end{equation}

\subsection*{2. Polynomial of degree 2}

Let
\begin{equation}
    P(t)=p_1t+p_2t^2,
    \qquad
    \mathbf{p}=\begin{bmatrix}p_1 & p_2\end{bmatrix}^{T}.
\end{equation}
Then
\begin{equation}
    P'(t)=p_1+2p_2t.
\end{equation}

\paragraph{(a) Quadratic cost matrix.}
Expanding and integrating the cost gives
\begin{align}
    J(P)
    &=\int_0^1 (p_1+2p_2t)^2\,dt \\
    &=\int_0^1
      \left(p_1^2+4p_1p_2t+4p_2^2t^2\right)dt \\
    &=p_1^2+2p_1p_2+\frac{4}{3}p_2^2.
\end{align}
Thus,
\begin{equation}
    J(P)=\mathbf{p}^{T}\mathbf{Q}\mathbf{p},
    \qquad
    \boxed{
    \mathbf{Q}=\begin{bmatrix}
        1 & 1 \\
        1 & \frac{4}{3}
    \end{bmatrix}.}
\end{equation}

\paragraph{(b) Equality constraint.}
The condition $P(1)=1$ is
\begin{equation}
    p_1+p_2=1.
\end{equation}
Therefore,
\begin{equation}
    \boxed{
    \mathbf{A}=\begin{bmatrix}1 & 1\end{bmatrix},
    \qquad
    \mathbf{b}=\begin{bmatrix}1\end{bmatrix}.}
\end{equation}

\paragraph{(c) Solution of the QP.}
From the equality constraint,
\begin{equation}
    p_1=1-p_2.
\end{equation}
Substituting this expression into the objective gives
\begin{align}
    J(p_2)
    &=(1-p_2)^2+2(1-p_2)p_2+\frac{4}{3}p_2^2 \\
    &=1+\frac{1}{3}p_2^2.
\end{align}
The unique minimizer is $p_2=0$, which implies $p_1=1$.
Consequently,
\begin{equation}
    \boxed{P^*(t)=t, \qquad J^*=1.}
\end{equation}
Allowing a quadratic term does not reduce the optimal cost.

\subsection*{3. Polynomial of degree 3}

Let
\begin{equation}
    P(t)=p_1t+p_2t^2+p_3t^3,
    \qquad
    \mathbf{p}=\begin{bmatrix}p_1 & p_2 & p_3\end{bmatrix}^{T}.
\end{equation}

\paragraph{(a) QP matrices.}
For coefficient indices $i,j\in\{1,2,3\}$,
\begin{equation}
    Q_{ij}
    =\int_0^1
      \left(it^{i-1}\right)
      \left(jt^{j-1}\right)dt
    =\frac{ij}{i+j-1}.
\end{equation}
It follows that
\begin{equation}
    \boxed{
    \mathbf{Q}=\begin{bmatrix}
        1 & 1 & 1 \\
        1 & \frac{4}{3} & \frac{3}{2} \\
        1 & \frac{3}{2} & \frac{9}{5}
    \end{bmatrix},
    \qquad
    \mathbf{A}=\begin{bmatrix}1 & 1 & 1\end{bmatrix},
    \qquad
    \mathbf{b}=\begin{bmatrix}1\end{bmatrix}.}
\end{equation}

\paragraph{(b) Solution of the QP.}
The Cauchy--Schwarz inequality gives a degree-independent lower bound:
\begin{align}
    J(P)
    &=\int_0^1 \left(P'(t)\right)^2dt \\
    &\geq
    \frac{\left(\int_0^1 P'(t)\,dt\right)^2}
         {\int_0^1 1^2\,dt} \\
    &=\left(P(1)-P(0)\right)^2 \\
    &=1.
\end{align}
Equality holds if and only if $P'(t)$ is constant. The endpoint
conditions then imply $P'(t)=1$, so
\begin{equation}
    p_1=1, \qquad p_2=0, \qquad p_3=0.
\end{equation}
Hence,
\begin{equation}
    \boxed{P^*(t)=t, \qquad J^*=1.}
\end{equation}
This result agrees with the Euler--Lagrange equation. When $r=1$,
the Euler--Lagrange equation reduces to
\begin{equation}
    P^{(2)}(t)=0,
\end{equation}
so the extremal is a polynomial of degree $2r-1=1$.

\subsection*{4. Additional terminal-velocity constraint}

Now consider
\begin{equation}
    \min_{P(t)} \int_0^1 \left(P'(t)\right)^2dt
\end{equation}
subject to
\begin{equation}
    P(0)=0, \qquad P(1)=1, \qquad P'(1)=-2.
\end{equation}

\paragraph{(a) Polynomial of degree 2.}
For
\begin{equation}
    P(t)=p_1t+p_2t^2,
\end{equation}
the constraints are
\begin{align}
    p_1+p_2&=1, \\
    p_1+2p_2&=-2.
\end{align}
Subtracting the first equation from the second gives $p_2=-3$, and
then $p_1=4$. Therefore,
\begin{equation}
    P^*(t)=4t-3t^2.
\end{equation}
The corresponding cost is
\begin{align}
    J^*
    &=4^2+2(4)(-3)+\frac{4}{3}(-3)^2 \\
    &=16-24+12 \\
    &=4.
\end{align}
Thus,
\begin{equation}
    \boxed{P^*(t)=4t-3t^2, \qquad J^*=4.}
\end{equation}

\paragraph{(b) Polynomial of degree 3.}
For
\begin{equation}
    P(t)=p_1t+p_2t^2+p_3t^3,
\end{equation}
the constraints are
\begin{align}
    p_1+p_2+p_3&=1, \\
    p_1+2p_2+3p_3&=-2.
\end{align}
Let $x=p_3$. Solving the constraints for the remaining coefficients gives
\begin{equation}
    p_2=-3-2x,
    \qquad
    p_1=4+x.
\end{equation}
Substituting these expressions into the quadratic cost yields
\begin{equation}
    J(x)=\frac{2}{15}x^2+x+4.
\end{equation}
Therefore,
\begin{equation}
    J'(x)=\frac{4}{15}x+1.
\end{equation}
Setting $J'(x)=0$ gives
\begin{equation}
    x=p_3=-\frac{15}{4}.
\end{equation}
Since $J''(x)=4/15>0$, this point is the unique minimizer. The other
coefficients are
\begin{equation}
    p_2=\frac{9}{2},
    \qquad
    p_1=\frac{1}{4}.
\end{equation}
Hence,
\begin{equation}
    P^*(t)
    =\frac{1}{4}t
     +\frac{9}{2}t^2
     -\frac{15}{4}t^3.
\end{equation}
The constraints are satisfied because
\begin{equation}
    P^*(0)=0,
    \qquad
    P^*(1)=1,
    \qquad
    {P^*}'(1)=-2.
\end{equation}
Finally,
\begin{align}
    J^*
    &=\frac{2}{15}\left(\frac{225}{16}\right)
      -\frac{15}{4}+4 \\
    &=\frac{17}{8}.
\end{align}
Therefore,
\begin{equation}
    \boxed{
    P^*(t)=\frac{1}{4}t+\frac{9}{2}t^2-\frac{15}{4}t^3,
    \qquad
    J^*=\frac{17}{8}.}
\end{equation}
The cubic polynomial achieves a lower cost than the quadratic polynomial:
\begin{equation}
    \frac{17}{8}<4.
\end{equation}

\section*{Deliverable 2: Multi-Segment Trajectory Optimization}

Let
\begin{equation}
    t_0<t_1<\cdots<t_k
\end{equation}
define $k$ trajectory segments. Segment $i$ is represented by
$P_i(t)$, and the prescribed scalar waypoints are
$\bar{x}_0,\bar{x}_1,\ldots,\bar{x}_k$. The following counts are for
one scalar coordinate.

\subsection*{1. Minimum snap trajectory: $r=4$}

For a minimum-snap trajectory,
\begin{equation}
    L=\left(P^{(4)}(t)\right)^2.
\end{equation}
The Euler--Lagrange equation reduces to
\begin{equation}
    P^{(8)}(t)=0.
\end{equation}
Therefore, each segment is a polynomial of degree $7$ and has $8$
coefficients. Across $k$ segments, there are $8k$ unknown coefficients.

\paragraph{Waypoint constraints.}
For each segment $i=1,\ldots,k$, impose
\begin{equation}
    P_i(t_{i-1})=\bar{x}_{i-1},
    \qquad
    P_i(t_i)=\bar{x}_i.
\end{equation}
There are two waypoint constraints per segment, for a total of
\begin{equation}
    2k
\end{equation}
waypoint constraints. Each internal waypoint is imposed once on the
segment arriving at the waypoint and once on the segment leaving it.

\paragraph{Free-derivative constraints.}
At every internal waypoint $t_i$, $i=1,\ldots,k-1$, impose derivative
continuity through order $6$:
\begin{equation}
    P_i^{(j)}(t_i)=P_{i+1}^{(j)}(t_i),
    \qquad
    j=1,\ldots,6.
\end{equation}
This gives
\begin{equation}
    6(k-1)
\end{equation}
free-derivative constraints.

\paragraph{Fixed-derivative constraints.}
At the initial and final times, prescribe velocity, acceleration, and jerk:
\begin{align}
    P_1^{(j)}(t_0)&=d_{0,j}, \\
    P_k^{(j)}(t_k)&=d_{k,j},
    \qquad j=1,2,3.
\end{align}
For a rest-to-rest trajectory, these prescribed values are normally zero.
There are three constraints at each endpoint, giving
\begin{equation}
    6
\end{equation}
fixed-derivative constraints.

The total number of constraints is therefore
\begin{equation}
    \boxed{
    2k+6(k-1)+6=8k,}
\end{equation}
which exactly matches the number of unknown polynomial coefficients.

\subsection*{2. General $r$-th derivative over $k$ segments}

For the general cost functional
\begin{equation}
    J(P)=\sum_{i=1}^{k}
    \int_{t_{i-1}}^{t_i}
    \left(P_i^{(r)}(t)\right)^2dt,
\end{equation}
the Euler--Lagrange equation on every segment is
\begin{equation}
    P_i^{(2r)}(t)=0.
\end{equation}
Thus, each segment has degree $2r-1$ and contains $2r$ coefficients.
The total number of unknown coefficients is
\begin{equation}
    2rk.
\end{equation}

The required constraints are as follows:
\begin{itemize}
    \item \textbf{Waypoint constraints:}
    two endpoint-position constraints per segment, giving
    \begin{equation}
        2k.
    \end{equation}

    \item \textbf{Free-derivative constraints:}
    at each of the $k-1$ internal waypoints, enforce continuity for
    derivative orders
    \begin{equation}
        j=1,\ldots,2r-2.
    \end{equation}
    This gives
    \begin{equation}
        (2r-2)(k-1)
    \end{equation}
    free-derivative constraints.

    \item \textbf{Fixed-derivative constraints:}
    at the initial and final times, prescribe derivative orders
    \begin{equation}
        j=1,\ldots,r-1.
    \end{equation}
    This gives
    \begin{equation}
        2(r-1)
    \end{equation}
    fixed-derivative constraints.
\end{itemize}

The total number of constraints is
\begin{align}
    2k+(2r-2)(k-1)+2(r-1)
    &=2k+(2r-2)k-(2r-2)+2r-2 \\
    &=2rk.
\end{align}
Therefore,
\begin{equation}
    \boxed{
    \begin{aligned}
        \text{polynomial degree per segment} &:\quad 2r-1, \\
        \text{waypoint constraints} &:\quad 2k, \\
        \text{free-derivative constraints} &:\quad (2r-2)(k-1), \\
        \text{fixed-derivative constraints} &:\quad 2(r-1).
    \end{aligned}}
\end{equation}

\end{document}
